% STATUS: written 2026-11-04, recalled results as is (one environment each, source kind and whole statement)
% Sources read in full: Martindale 1964 (pp. 183--187).
\section{The problem of Martindale and the standard form}\label{sec:lie-type:martindale}

Chapter~\ref{ch:jordan-type} followed the Jordan product $a\circ b=ab+ba$. We now follow
the other product that an associative ring carries, the commutator $[a,b]=ab-ba$. A map
that satisfies the product rule for the commutator is a Lie derivation \cite[p.~183]{Martindale1964}.
Martindale began the theory with the definition \cite[p.~183]{Martindale1964}.

\begin{definition}[{\cite[p.~183]{Martindale1964}}]\label{def:lie:martindale}
Let $A$ be an associative ring. An additive mapping $a\mapsto a'$ of $A$ into itself is
called a \emph{Lie derivation} of $A$ if for every pair $a,b\in A$,
$[a,b]'=[a',b]+[a,b']$.
\end{definition}

Every derivation is a Lie derivation. % claim-ok: a trivial check
So is the sum of a derivation and an additive map into the center that sends
commutators to zero. The problem, set by Martindale in his first paragraph, is
whether these two pieces exhaust the Lie derivations \cite[p.~183]{Martindale1964}.
Herstein had studied Lie derivations of the classical noncommutative simple rings; Martindale
extends the theory to rings whose nonzero two-sided ideals have a nonzero right
annihilator, and he notes that primitive rings are precisely such rings
\cite[p.~183]{Martindale1964}.

His first result extends the map itself. % claim-ok: narration of the chapter plan, no mathematical content

\begin{theorem}[{\cite[Theorem~1, p.~186]{Martindale1964}}]\label{thm:lie:martindale-extension}
Let $A$ be a primitive ring containing a nontrivial idempotent, but not a domain, with
characteristic different from $2$. Then every Lie derivation of $A$ can be uniquely
extended to an ordinary derivation of an associative algebra $A^{\ast}$ containing $A$.
\end{theorem}

The theorem is stated for a primitive ring $A$ that contains a nontrivial
idempotent, in characteristic not two \cite[pp.~183--184]{Martindale1964}. That
is the setting the decomposition needs, and the paper states it there.
The ring of Martindale quotients is a different construction, built in 1969
\cite[p.~576]{Martindale1969}; Chapter~\ref{ch:prime-semiprime} uses it for
semiprime rings, and this theorem does not need it. % claim-ok: narration of the chapter plan, no mathematical content

\begin{theorem}[{\cite[Theorem~2, p.~187]{Martindale1964}}]\label{thm:lie:martindale-standard}
Let $A$ and $A^{\ast}$ be as in Theorem~\ref{thm:lie:martindale-extension}. Every Lie
derivation $d$ of $A$ is of the form $d=d'+\tau$, where $d'$ is the unique extension
of $d$ to an ordinary derivation of $A^{\ast}$, and $\tau$ is an additive mapping
of $A$ into the center of $A^{\ast}$. The map $\tau$ sends commutators into zero.
\end{theorem}

The form $d=\text{derivation}+\text{central map vanishing on commutators}$ is the
answer the chapter keeps meeting. % claim-ok: the recalled theorems just above state this form
The functional identities of Section~\ref{sec:lie-type:bresar} later settle when the
central map really occurs.
Martindale closes with simple rings, where the map into the center stays inside
the ring \cite[p.~187]{Martindale1964}.

\begin{remark}[{\cite[p.~187]{Martindale1964}}]\label{rem:lie:martindale-simple}
Let $A$ be a simple ring characteristic not $2$ with minimal one-sided ideals. Then
every Lie derivation of $A$ is of the form $d+\tau$. Here $d$ is an ordinary
derivation of $A$ into itself, and $\tau$ is additive. The map $\tau$ sends $A$
into the center of $A$ and commutators into zero. % claim-ok: the source remark quoted in this environment states $\tau$ maps $A$ into the center and commutators to zero
\end{remark}

Remark~\ref{rem:lie:martindale-simple} settles the simple case that has minimal
one-sided ideals, and the general case is the problem of Herstein, still open. % claim-ok: own signpost introducing the moved problem, whose statement carries its own citation

\begin{problem}[{\cite[Problem~3, p.~529]{Herstein1961Lie}}]\label{prob:open:herstein}
Let $R$ be a simple ring.
Is every Lie derivation of $R$ an ordinary derivation plus an additive map
into the center?
\end{problem}

Herstein raised the question for a simple ring, and he recorded that Kaplansky
had settled the case of matrix units of size at least $3$
\cite[Problem~3, p.~529]{Herstein1961Lie}. Martindale proved his theorem for
primitive rings with a nontrivial idempotent and said that it falls far short
of a general solution of that conjecture \cite[p.~183]{Martindale1964}.
Theorem~\ref{thm:lie:martindale-standard} gives the form for a primitive ring
with a nontrivial idempotent \cite[Theorem~2, p.~187]{Martindale1964}. His
closing remark covers the simple case: the derivation then maps $R$ into
itself and the additive map goes into the center of $R$
\cite[p.~187]{Martindale1964}.

Bre\v{s}ar extends the form to every prime ring of characteristic not $2$ that
does not satisfy the identity $S_4$
\cite[Theorem~5, p.~529]{Bresar1993LieTAMS}. His Lemma~1 lists the rings that
do satisfy $S_4$: they are the commutative rings and the subrings of
$M_2(F)$ for a field $F$ \cite[Lemma~1, p.~530]{Bresar1993LieTAMS}. He records
that the assumption on $S_4$ may be removable \cite[p.~528]{Bresar1993LieTAMS}.
So Problem~\ref{prob:open:herstein} stays open in characteristic $2$ and for the
rings that satisfy $S_4$.

\hbfigure{05-lie-type/standard-decomposition}{The standard form of Martindale \cite[Theorem~2, p.~187]{Martindale1964}. A Lie derivation $d$ splits as $d=d'+\tau$, where $d'$ is an ordinary derivation of $A^{\ast}$ and $\tau$ is an additive map into the center of $A^{\ast}$ that sends commutators to zero. Theorem~\ref{thm:lie:martindale-standard} states it in full.}{fig:lie:standard-decomposition}
